%\needspace{4\baselineskip}
\color{uniblue}{\subsection{Логика сигнализации}}
\color{black}

\begin{enumerate}[label=\arabic{section}.\arabic{subsection}.\arabic{enumi}, labelsep=4pt, leftmargin=0pt, itemindent=57pt, itemsep=0pt, parsep=0pt] % если что можно крутить itemsep


\item
В логике сигнализации контролируются уровни изоляции ТТ. Накладка «Откл. авар. Изол. ТТ» 
позволяет сформировать сигнал на отключеие.

\item
Устройство формирует выходное воздействие на сигнализацию «Сигн. сраб.» и «Сигн. сраб. 
(фикс.)» по схеме «ИЛИ» в следующих случаях: 

\begin{itemize}
    \item Отключении БСК; 
    \item Срабатывании сигнальной ступени ЗВП;
    \item Срабатывании сигнальной ступени ЗПВГ; 
    \item Срабатывании сигнальной ступени ЗПН; 
    \item Срабатывании ЗНР; 
    \item Срабатывании УРОВ.
\end{itemize}

\item
Устройство формирует выходное воздействие на сигнализацию «Сигн. неиспр.» и «Сигн. неиспр. 
(фикс.)» по схеме «ИЛИ» в следующих случаях:

\begin{itemize}
    \item Неиспр. ЦТ ДТЗ; 
    \item Неиспр. ЦТ ДТЗ НП;
    \item Неиспр. ТТ; 
    \item Неиспр. ОТ терм;
    \item Неиспр. ЦН.  
\end{itemize}

\item
Сброс фиксации производится с помощь «Сброс» 


\begin{figure}[H]
\centering
\includegraphics[width=0.75\textwidth,height=0.75\textheight,keepaspectratio]{\fbpath/91 Логика сигнализации/img/Log_sign.pdf}
\captionof{figure}{Функциональный блок логики сигнализации}\label{pict:lo_sign}
\end{figure}


\small
\setlength{\extrarowheight}{0.05cm} % высота ячеек в таблице
\begin{longtable}
{|>{\centering\arraybackslash}m{4.5cm}|>{\centering\arraybackslash}m{4.5cm}|>{\centering\arraybackslash}m{3cm}|>{\centering\arraybackslash}m{3.5cm}|}
\caption{Перечень уставок БНН\hfill\vspace{-0.5\baselineskip}}\label{ksn:tbl2} \\ 
\hhline{|-|-|-|-|}
\rowcolor{gray!30}
Наименование уставки & Диапазон & Шаг &  Значение по умолчанию \\ 
\hhline{|-|-|-|-|}
\endfirsthead

\caption*{\hspace{3pt}\emph{Продолжение таблицы \ref{ksn:tbl2}\hfill\vspace{-0.5\baselineskip}}} \\ % сделано по ГОСТ 2.105 п.6.8.7
\hhline{|-|-|-|-|}
\rowcolor{gray!30}
Наименование уставки & Диапазон & Шаг &  Значение по умолчанию \\  
\hhline{|-|-|-|-|}
\endhead
\hhline{|-|-|-|-|}

\endfoot
\endlastfoot

%%%%%%%%%%%%содержимое%%%%%%%%%%%%%%

\end{longtable}
\setlength{\extrarowheight}{0.0cm} % возврат параметра высота ячеек в таблице
\normalsize




















\end{enumerate} 